% New English translation from checked cote 25, site 57 (PDF page 58).
% AI-translated by GPT-6 Astra (OpenAI, ChatGPT Pro mode).
\scanpage{25/57}
\mspage{\unsure{52|92}}

The hypothesis dim $f(X)\hand{\geqslant} 2$ means that the transcendence degree of
$k(t_1,..,t_r)$ over $k$ is $\hand{\geqslant} 2$, i.e. we can find $t_1,t_2$, say, alg independent.
In a), we then take $P = T_1$, $Q = T_1^{\,2}$, $R = T_1T_2$, In b), we
do the same, \var{\xout{but taking care to take $t_1$}}{}{typewriter strike-out}{ega25_25_057-v0001} noting that above we can
choose $t_1$ coming from a $T_1$ \uline{vanishing on} L, thanks to the fact $f(X)\hand{\not\subset} L$
(which implies that there exists an index i between 1 and s such that $t_s$ is
nonzero, hence $t_s$ is nonconstant (since $t_s$ vanishes at x), hence $t_s$ is not alg over k
(NB we may assume k alg closed). Case c) follows from b), \var{\deleted{even with n=2
instead of n=3,}}{}{deletion}{ega25_25_057-v0002} except in the case where \var{\xout{X}}{}{typewriter strike-out}{ega25_25_057-v0003} f(X) is contained in
the image \var{\inserted{L}}{L}{insertion}{ega25_25_057-v0004} under \var{\over{\unsure{f}}{f}}{f}{writing over}{ega25_25_057-v0005} of the tangent space of X at x, \var{\deleted{which moreover requires for
dimension reasons that the tangent map at x be injective, i.e.
f unramified at x, \var{\xout{When dim X}}{}{typewriter strike-out}{ega25_25_057-v0007} so that f(X) is contained in
a linear subvariety of P of the same dimension as X.}}{}{deletion}{ega25_25_057-v0006} \hand{[}When dim X=2,
this case is indeed exceptional (the trace on a plane of a quadric tangent
to that plane generally consists of two intersecting
lines, and hence is not irreducible)\hand{]}. But to treat this case, \var{\xout{under}}{}{typewriter strike-out}{ega25_25_057-v0008} the
form P,Q,R made explicit above, we can clearly replace X by
L itself, where the solution is trivial (\var{\inserted{if dim L $=2$}}{if dim L $=2$}{insertion}{ega25_25_057-v0009} take $P=T_r^{\,2}$\var{\xout{\illegible{1}}}{}{typewriter strike-out}{ega25_25_057-v0010}, $Q=T_r^{\,3}$,
$R = T_r^{\,2}T_{r-1}$, noting that $\unsure{Q}/\unsure{P} = T_r$ and $R/P = T_{r-1}$ are
linear forms lin independent on L, hence alg independent\var{\inserted{if dim $L\geqslant 3$, then $T_{r-2}$, $T_{r-1}$, $T_r$ are lin. independent on L, and we take $P=T_r^{\,2}$, $Q=T_{r-1}^{\,2}$, $R=T_{r-2}^{\,2}$.}}{ if dim $L\geqslant 3$, then $T_{r-2}$, $T_{r-1}$, $T_r$ are lin. independent on L, and we take $P=T_r^{\,2}$, $Q=T_{r-1}^{\,2}$, $R=T_{r-2}^{\,2}$.}{insertion}{ega25_25_057-v0011}). \var{\deleted{We could mention in a
remark that in c) we can even take n=2 unless f(X) is contained
in a plane.}}{}{deletion}{ega25_25_057-v0012}

Combining this with \unsure{8.181|8.18L|8.18}, we find a

\uline{Corollary} 9.4. \var{\xout{(For X proper over k}}{}{typewriter strike-out}{ega25_25_057-v0013}

Finally, we must also combine the latter with 9.2 to
find a theorem summing up the “excellent case”:
\var{\over{\uline{\unsure{Corollary}}}{\uline{Theorem}}}{\uline{Theorem}}{writing over}{ega25_25_057-v0014} 9.5. (When X is smooth \var{\deleted{and}}{}{deletion}{ega25_25_057-v0015} proper\var{\inserted{(and geom integral)}}{ (and geom integral)}{insertion}{ega25_25_057-v0016} over k, and $f:X\hand{\to} P$ unramified,
\hand{X}of dimension $\hand{\geqslant} 2$)\hand{, mentioning the result on \unsure{$\pi_1$} ?.}
\hand{when $X\to P$ \unsure{not|an} \unsure{imm|immer} (variety of singular pts of \unsure{$f$}).}
\tnote{25-057-1}{The repeated subscript s after the choice of an index i is retained. The conclusion of Corollary9.4 is not supplied, and the unfinished formulation of Theorem9.5 and its handwritten notes remain unfinished. The canceled sous leaves a fragmentary construction before the form P,Q,R; it is retained.}

\var{\margin{57}}{\placedmargin{57}}{marginal note placed}{ega25_25_057-v0017}
